\section{Methods}

\subsection{Quantum Prisoner's Dilemma}

The following construction summarizes the standard EWL protocol \cite{ewl}. The classical actions cooperate and defect are represented by the computational basis states

\begin{equation}
	|C\rangle,\quad |D\rangle.
\end{equation}

The initial state is

\begin{equation}
	|CC\rangle.
\end{equation}

An entangling operator $J$ is applied to create quantum correlations.

Each player then applies a local unitary strategy:

\begin{equation}
	U_A \otimes U_B.
\end{equation}

Finally, the entanglement is reversed using $J^\dagger$ and the state is measured.

The final state is

\begin{equation}
	|\psi_f\rangle
	=
	J^\dagger
	(U_A \otimes U_B)
	J
	|CC\rangle.
\end{equation}

After measurement, the final state produces probabilities

\begin{equation}
	P_{CC},\quad P_{CD},\quad P_{DC},\quad P_{DD}.
\end{equation}

Expected payoffs are computed as

\begin{equation}
	E[u_A]
	=
	RP_{CC}
	+
	SP_{CD}
	+
	TP_{DC}
	+
	PP_{DD},
\end{equation}

and

\begin{equation}
	E[u_B]
	=
	RP_{CC}
	+
	TP_{CD}
	+
	SP_{DC}
	+
	PP_{DD}.
\end{equation}

\subsection{Entanglement Operator}

Let $D$ denote the matrix representation of the classical defection operator. The entangling operator is commonly defined as
\begin{equation}
	J(\gamma)=\exp\left(
	-\frac{i\gamma}{2}D\otimes D
	\right),
\end{equation}

where

\begin{equation}
	\gamma \in \left[0,\frac{\pi}{2}\right].
\end{equation}

Unlike classical mixed-strategy games, the resulting payoff functions depend continuously on unitary parameters and quantum correlations induced by entanglement. This continuous structure motivates interpreting equilibrium search as optimization over continuous parameterized unitary strategy spaces.

The parameter $\gamma$ controls the degree of entanglement within the EWL framework. Specifically, $\gamma = 0$ corresponds to the classical game, whereas $\gamma = \pi/2$ corresponds to maximal entanglement.

\subsection{Quantum Strategies}

Players choose unitary strategies $U(\theta,\phi)\in SU(2)$ parameterized by angular variables $(\theta,\phi)$:
\begin{equation}
	U(\theta,\phi)
	=
	\begin{pmatrix}
		e^{i\phi}\cos(\theta/2)
		&
		\sin(\theta/2)
		\\
		-\sin(\theta/2)
		&
		e^{-i\phi}\cos(\theta/2)
	\end{pmatrix}.
\end{equation}

Classical cooperation and defection appear as special cases:

\begin{equation}
	C = I,
\end{equation}

\begin{equation}
	D =
	\begin{pmatrix}
		0 & 1 \\
		-1 & 0
	\end{pmatrix}.
\end{equation}

EWL also defines the quantum strategy

\begin{equation}
	Q =
	\begin{pmatrix}
		i & 0 \\
		0 & -i
	\end{pmatrix}.
\end{equation}

\subsection{Equilibrium Search as Optimization}

For fixed opponent strategy $U_B$, player $A$ solves

\begin{equation}
	U_A^*
	\in
	\arg\max_{U_A \in \mathcal{S}}
	E[u_A(U_A,U_B;\gamma)].
\end{equation}

Similarly, player $B$ solves

\begin{equation}
	U_B^*
	\in
	\arg\max_{U_B \in \mathcal{S}}
	E[u_B(U_A,U_B;\gamma)].
\end{equation}

A Nash equilibrium is therefore a fixed point of coupled optimization problems.

Unlike classical finite games, the quantum strategy space is continuous. Because admissible quantum strategies are parameterized continuously, equilibrium computation may be interpreted as a coupled nonconvex optimization problem over continuous unitary strategy spaces. Let each player parameterize a strategy by angular variables $(\theta,\phi)$ and iteratively update these parameters through gradient ascent on the expected payoff functional:

\begin{equation}
	(\theta_A,\phi_A)
	\leftarrow
	(\theta_A,\phi_A)
	+
	\eta \nabla E[u_A],
\end{equation}

\begin{equation}
	(\theta_B,\phi_B)
	\leftarrow
	(\theta_B,\phi_B)
	+
	\eta \nabla E[u_B].
\end{equation}

Here, $\eta > 0$ denotes a learning-rate parameter, while $\nabla E[u_i]$ represents the gradient of the expected payoff with respect to the player's strategy parameters. These updates are intended as conceptual analogies to policy-gradient learning rather than convergence-guaranteed optimization procedures.

This optimization perspective connects quantum equilibrium search to reinforcement learning and continuous policy optimization \cite{sutton,busoniu}.

From the perspective of statistical learning, expected payoffs may also be interpreted as stochastic objective functions generated by probabilistic measurement outcomes. Consequently, equilibrium search in quantum strategic games may be interpreted as optimization over continuous probabilistic policy spaces rather than equilibrium selection over finite discrete action spaces.